\documentclass[11pt]{article}
\usepackage[T1]{fontenc}
\usepackage{lmodern}
\usepackage[a4paper,margin=30mm]{geometry}
\usepackage{amsmath,amssymb,amsthm,mathtools}
\usepackage{microtype}
\usepackage[numbers,sort&compress]{natbib}
\usepackage{booktabs}
\usepackage{xurl}
\usepackage{xcolor}
\definecolor{linkblue}{RGB}{20,54,100}
\usepackage[colorlinks=true,allcolors=linkblue]{hyperref}
\hypersetup{pdftitle={Information, Incentives and Symmetric Equilibria in Finite Coordination Games},pdfauthor={Nicholas Ham},pdfsubject={Finite coordination, information and verification},pdfkeywords={coordination, equilibrium, verification, symmetry}}
\providecommand{\doi}[1]{\href{https://doi.org/#1}{\nolinkurl{#1}}}
\newtheorem{theorem}{Theorem}[section]
\newtheorem{lemma}[theorem]{Lemma}
\newtheorem{proposition}[theorem]{Proposition}
\newtheorem{corollary}[theorem]{Corollary}
\theoremstyle{definition}
\newtheorem{definition}[theorem]{Definition}
\newtheorem{example}[theorem]{Example}
\theoremstyle{remark}
\newtheorem{remark}[theorem]{Remark}
\newtheorem{question}[theorem]{Open question}
\DeclareMathOperator{\Aut}{Aut}
\DeclareMathOperator{\Sym}{Sym}
\newcommand{\R}{\mathbb R}
\newcommand{\Sn}{\mathfrak S_n}
\newcommand{\id}{\mathrm{id}}

\title{Information, Incentives and Symmetric Equilibria\\in Finite Coordination Games}
\author{Nicholas Ham}
\date{}
\begin{document}
\maketitle
\begin{abstract}
We study finite coordination games through three separate constraints:
information, incentives and representation. A guaranteed-success screening
bound identifies the sharp no-transfer welfare frontier for independent
private costs and general upward-closed contribution tasks. Trusted binary
cost verification gives an exact arbitrary-player welfare--query frontier.
For explicitly tabulated local-factor games, an action-orbit polynomial
reduction yields parameterized invariant Nash computation, conditional on
standard real-algebraic sampling bounds. A one-action-orbit circuit game
separates easy equilibrium output from hard exact payoff evaluation.
Symmetry averaging preserves convex mechanism objectives but can change
implementation resources. These are specializations and adaptations of
established screening, search and equilibrium methods; originality remains
provisional. The earlier allocation and state models are retained as context.
No empirical claim about language-model rationality is made.
\end{abstract}
\noindent\textbf{Keywords:} finite games; coordination; private information;
verification; symmetric equilibria.

\section{Models, observations and coordination}\label{sec:01-introduction-1}
A coordinator can make the right decision without identifying a whole utility
table; identifying a relabeling need not reveal enough to optimize a task.
Similarly, a symmetric mixed equilibrium can exist while an efficient
realization requires a coordination resource. We examine these distinctions
in finite models with explicit timing, information and receipts.

Other-play and label-free coordination already investigate symmetry and
independently constructed policies \citep{hu2020,treutlein2021}. Neural
solvers and equivariant approximators directly study equilibrium selection
and welfare limitations \citep{marris2022,duan2023}. Tewolde et al.\
\citep{tewolde2025} connect game automorphisms with computation and recall
the classical existence of invariant mixed equilibria. The contribution
assessed here is therefore a set of precise model comparisons, with novelty
unresolved, rather than a new general symmetry or equilibrium theory.

The mathematical comparisons concern finite coordination models with explicit
information and implementation resources. The action-orbit computation results
use typed game automorphisms; the utility comparisons retain a specified common
receipt scale. Full structural classification and general assignment fairness
are separate questions and are not prerequisites for the results here.

All initial games and state sets are finite. A typed relabeling moves players
and bijects their action sets; products act right to left. Exact maps preserve
numerical utilities, cardinal maps preserve each recipient's expected-utility
preferences, and ordinal maps preserve pure-outcome preferences. Interpersonal
sums below use explicitly calibrated receipts. For software agents such
objectives are design specifications, not independently established welfare
interests. We distinguish pure equilibria, independent mixed profiles,
lotteries over equilibria, and private recommendations.

Sections~\ref{sec:02-information-1}--\ref{sec:05-state-1} give the allocation,
equilibrium-selection, private-cost and state models. The main refinements are
the guaranteed-success screening bound
(Proposition~\ref{result:07-phase11-1}), the sharp trusted-query frontier
(Proposition~\ref{result:07-phase11-2}), and the explicitly tabulated local-factor
reduction (Proposition~\ref{result:07-phase11-3}). The representation boundary
and implementation costs explain why these conclusions require different
resources. Each result retains its finite model and stated assumptions; no
claim about general-purpose language-model behaviour follows.

\section{Information sufficient for allocation}
A coordinator observes $Lw$ before choosing an assignment from a nonempty
finite set $D$. The unknown capability vector belongs to an unrestricted
finite-dimensional real vector space $W$; assignment $d$ gives disclosed
linear task reward $b_d(w)$. Observations are truthful, fixed linear
measurements. No strategic reporting, prior, positivity restriction or
computational bound is assumed. Fix $d_0$ and let $B$ have rows $b_d-b_{d_0}$.

\begin{theorem}[Allocation contrasts; AG-A1]
An $Lw$-based deterministic rule can choose an optimal assignment for every
$w\in W$ iff $\ker L\subseteq\ker B$. The minimum unrestricted scalar
linear measurement rank is $\operatorname{rank}B$.
\end{theorem}
\begin{proof}
For $t\in\ker L$, the assignment chosen at zero data must maximize reward
at both $t$ and $-t$. Every reward functional consequently agrees on $t$,
so $Bt=0$. Conversely the kernel inclusion makes all reward contrasts
well-defined functions of $Lw$, which suffice for optimization. Rank-nullity
gives the lower bound; measuring $Bw$ attains it.
\end{proof}
This is a standard linear-estimability argument applied to decisions, not a
new general theory of decision-focused prediction \citep{elmachtoub2020}.
For a restricted capability domain $C$, the correct condition is a common
maximizer on each observation fiber:
\[
 \bigcap_{w\in C:\,Lw=y}\arg\max_{d\in D}b_d(w)\ne\varnothing.
\]
For example, $b_0=0,b_1(w)=w$ on $w\ge0$ needs no measurements, unlike
the full-line problem. Offsets and bounded domains cannot be ignored.

For $n$ agents filling $n$ fixed tasks bijectively, write $C_{ij}$ for arbitrary
real completion scores. Assignment $\pi$ has score $\sum_i C_{i,\pi(i)}$.
The reward-contrast span consists of matrices with zero row and column sums,
with basis
$E_{ij}-E_{in}-E_{nj}+E_{nn}$, $i,j<n$. Each basis rectangle is a difference
of two assignments exchanging two task choices. Thus exactly $(n-1)^2$
unrestricted scalar linear measurements suffice and are necessary.

By comparison, testing physical capability-row interchangeability of two agents
with tasks fixed has discrepancy $T(C)=(C_{1j}-C_{2j})_j$, rank $n$.
This is a capability-array test, not a claim about a complete normal-form game's
automorphism group. The canonical Game Fairness criterion identifies $Tw=0$
through linear observations exactly when their kernel is contained in $\ker T$.
For $n=2$ allocation rank $1$ is smaller than test rank $2$; for $n=3$
allocation rank $4$ is larger than test rank $3$. Their row spaces at $n=3$
are incomparable: a constant added to one row changes interchangeability but
not any assignment comparison; changing only $C_{31}$ from positive to
negative leaves the first two rows equal but changes which tasks can be
optimal. The joint measurement rank is $5$, since the row-space intersection
has dimension $2$.

\paragraph{Noise, messages and incentives.}
If reward contrasts factor as $b_d-b_{d_0}=a_dL$ and observations have error
$\|e\|\le\epsilon$, maximizing $a_d\cdot(Lw+e)$ incurs regret at most
\[
 \epsilon\max_{d,e}\|a_d-a_e\|_*.
\]
Observed optimality and the dual-norm inequality prove this; for two decisions
it is the sharp deterministic minimax value on a full linear class with
surjective $L$. It is not a randomized-regret or finite-bit bound. An informed
sender can instead communicate a selected assignment using $n!$ messages,
which is sharp since each permutation is uniquely optimal somewhere. Nor does
truthful measurement imply truthful reports: a low-capability agent paid only
for receiving a task may profitably exaggerate its capability. Full proofs
and exact witnesses remain in the allocation research note.

\section{Symmetric selection and implementation resources}
The suggested obstruction ``automorphisms move every equilibrium'' must be
qualified: it cannot exclude a fixed independent mixed Nash profile in a
finite game with cardinal automorphisms.

\begin{proposition}[Classical invariant equilibrium; AG-E1]
For every finite group $G$ of cardinal typed automorphisms, a finite game
has a mixed Nash equilibrium fixed by every $g\in G$.
\end{proposition}
\begin{proof}
The fixed set $X^G$ in the product of mixed-strategy simplexes is nonempty,
compact and convex, by affine averaging. At fixed $x$, cardinal maps preserve
best responses, so $\mathrm{BR}(x)$ is $G$-stable. Averaging a best-response
profile stays in its convex best-response product and is fixed by $G$.
Thus $\mathrm{BR}(x)\cap X^G$ is nonempty compact convex with closed graph.
Kakutani gives a fixed point, hence the required equilibrium.
\end{proof}
For exact maps this is Nash's established theorem, recalled by
\citet[Theorem 7, long version]{tewolde2025}. Cardinal best-response preservation
is prior Game Fairness U5. Ordinal automorphisms need not preserve mixed
best responses and cannot be substituted. We have not proved existence by
averaging arbitrary equilibria: their marginal average can fail to be Nash.

For a game class and relabeling action, a set-theoretic equivariant selector
from a relabeling-equivariant solution correspondence exists iff every input stabilizer fixes some
solution. Choose a fixed solution per input orbit and transport it to prove
sufficiency; necessity is immediate. This supplies no continuity, computation
or decentralized implementation guarantee. Averaging a point mass over a finite
group always supplies an invariant lottery on a nonempty invariant solution
set, even if that set is infinite. If its support consists of Nash profiles,
its induced action distribution is correlated equilibrium by linearity of
obedience inequalities. It is not generally a product distribution.

For example, two agents each receive $1$ for unequal binary actions and $0$
otherwise. Pure Nash profiles $(0,1),(1,0)$ are exchanged by the exact player
swap. Neither is fixed, while the unique fully mixed equilibrium is the fixed
profile $(1/2,1/2)$, with task success $1/2$. An invariant lottery on the two
pure profiles attains success $1$. This classical anti-coordination witness
already appears in \citet{cheng2004}; the broader welfare obstruction for
equivariant approximators has substantial prior treatment \citep{duan2023}.

\begin{proposition}[An anonymous resource bound; AG-E2]
Suppose the agents use identical local code and observations, public randomness
with the same interpretation, and conditionally IID private random
seeds. Without distinguishing roles or communication, anti-coordination
success is at most $1/2$. A private complementary recommendation device
attains $1$ obediently.
\end{proposition}
\begin{proof}
Conditional on a public seed, actions are IID Bernoulli$(p)$ and success is
$2p(1-p)\le1/2$. Average over seeds. Complementary recommendations always
succeed, and no recommended player gains by changing its action.
\end{proof}
Distinguishable roles with a public fair bit also attain $1$; they drop the
identical-observation assumption. An invariant distribution existing
mathematically is therefore separate from its available implementation.

\section{Private costs and an exact mediation gap}\label{sec:04-bayesian-1}
Two agents independently draw private costs $\theta_i\in\{1/2,3/2\}$ with equal
probabilities. Each observes its own cost before choosing whether to volunteer,
$a_i\in\{0,1\}$, simultaneously. Their common prior and utility model are known:
\[
 u_i(a,\theta)=2\mathbf1_{a_0+a_1\ge1}-\theta_i a_i,\qquad
 W=4\mathbf1_{a_0+a_1\ge1}-\theta_0a_0-\theta_1a_1.
\]
There are initially no reports, payments or private correlated recommendations.
The environment is exchange-symmetric ex ante; a mixed-cost realized game is
not symmetric under bare player exchange. Interim beliefs follow the specified
prior, rather than an assumed complete-information theorem.

\begin{theorem}[Private information and architecture]\label{result:04-bayesian-1}
The unique symmetric behavioral Bayesian Nash equilibrium volunteers at low
cost only, with welfare $5/2$. Among all possibly distinct local policies with
public randomness independent of types and independent private seeds, welfare
is at most $3$, sharply. An externally informed private mediator can attain
the ex-post optimum, welfare $13/4$ in expectation, by an exchange-equivariant
obedient recommendation scheme.
\end{theorem}
\begin{proof}
Let $\alpha,\beta$ be the common probabilities at low and high cost, and
$q=(\alpha+\beta)/2$. Volunteering's interim gain is $2(1-q)-\theta$.
If $q<1/4$ both types must volunteer; if $q>3/4$ both must abstain, each a
contradiction. At $q=1/4$ low must volunteer, forcing $q\ge1/2$; at $q=3/4$
high must abstain, forcing $q\le1/2$. Otherwise low volunteers and high
abstains, giving $q=1/2$ and strict gains $1/2,-1/2$. Direct averaging gives
welfare $5/2$.

Conditional on a public seed, welfare is multilinear in the four local
behavioral probabilities, so a maximum occurs at deterministic policies.
With $N,L,H,A$ denoting never, low only, high only and always volunteering,
all possibilities have the welfare values in Table~\ref{tab:volunteer-welfare}.
\begin{table}[htbp]
\centering
\caption{Welfare by deterministic local volunteering policies.}\label{tab:volunteer-welfare}
\begin{tabular}{crrrr}
\toprule
 &N&L&H&A\\
\midrule
N&0&7/4&5/4&3\\
L&7/4&5/2&2&11/4\\
H&5/4&2&3/2&9/4\\
A&3&11/4&9/4&2\\
\bottomrule
\end{tabular}
\end{table}
The largest entry is $3$. Public randomization cannot improve the bound.
A designated volunteer attains it and is a Bayesian equilibrium; a public
fair lottery between the two designated equilibria makes the selection law
exchange-invariant.

Finally let a trusted mediator observe actual costs and recommend exactly one
minimum-cost volunteer, uniformly at ties. At every realized type profile,
the volunteer gains $2-\theta_i>0$ from obeying; the abstainer would only
pay an extra cost by deviating. Hence the scheme obeys conditional incentives
as well. One cheapest volunteer maximizes $W$ pointwise, and its expected
cost is $3/4$, giving $13/4$. The rule commutes with exchange in distribution.
\end{proof}
This is a finite Bayesian-correlation comparison in the established setting
of \citet{bergemann2016}, not a new general mediation theorem. External type
observation is substantive. A low type can profitably report high if reports
are trusted: it reduces its selection probability while another volunteer
still supplies the benefit. Obedience is not truthful elicitation.

The bounds also depend on implementation. Requiring identical local policies
lowers the best unconstrained welfare to $41/16$, attained at
$(\alpha,\beta)=(1,1/4)$, which is not Bayesian Nash. Indeed with
$q=(\alpha+\beta)/2$ welfare is
$8q-4q^2-\alpha/2-3\beta/2$; minimizing costs at fixed $q$ volunteers low
first and yields this maximum by completing the square. Thus the symmetric
BNE value $5/2$, optimal anonymous-policy value $41/16$, public labeled-policy
value $3$ and informed mediator value $13/4$ concern different feasible sets.

Invariant Bayesian equilibrium existence has a separate standard proof: with
positive marginal type probabilities, take the product of every type's
behavioral simplex. For a finite group preserving prior, typed actions and
\emph{interim} best responses, the fixed-space Kakutani argument applies to
the conditional BR product. Ex-ante numerical symmetry alone is not a license
to omit those assumptions. This existence does not supply the mediator's
information or efficiency.

\section{Persistent state and priced checkpoint tracking}\label{sec:05-state-1}
Observable path independence after invisible transformations are quotiented
out is an algebraic condition. The following operational model instead fixes
actuators, targets and prices to distinguish splitting from tracking cost.

A two-role carrier has state $x\in\mathbb Z/4$, role parity $p=x\bmod2$ and
calibration bit $b=\lfloor x/2\rfloor$. The task succeeds iff $b=0$, so its
accepted configurations are $0,1$. Starting at $0$, it receives $T$ swap
requests, each followed by a task checkpoint. The required state is the
current parity, not just any state with that parity. The fixed primary
actuator adds $1$; after each actuation a controller may add $2$ at cost $c>0$.
One primary actuation per request is mandatory, and extra primary operations
are unavailable. These targets and prices are physical model inputs.

\begin{proposition}[Priced checkpoint tracking]\label{result:05-state-1}
The minimum correction cost for success at every checkpoint is
$c\lfloor T/2\rfloor$. If both primary directions $+1,-1$ are freely
available and the controller can observe or remember current parity, every
checkpoint succeeds with zero correction cost, although
$\mathbb Z/4\to\mathbb Z/2$ does not split.
\end{proposition}
\begin{proof}
Successful states alternate $0,1$. From $0$, adding $1$ reaches $1$ for free;
from $1$ it reaches $2$ and needs one costed addition of $2$ to reach $0$.
These forced checkpoint states prove necessity and attainability. With both
actuators, choose $+1$ at $0$ and $-1$ at $1$. The odd elements nevertheless
have square $2$, preventing a homomorphic section.
\end{proof}
If inverse actuation costs $d\ge0$, the minimum is
$\lfloor T/2\rfloor\min(c,d)$. When reverse actuation is free, a controller with no clock, feedback or other
input than identical swap requests needs two internal states for zero-cost
tracking with both actuators; one deterministic state repeats a fixed direction
and fails one of the first two checkpoints. If parity is observed, that memory
is already supplied. A one-bit reset-at-each-call message is a different
resource from persistent local memory and is not needed per request here.

More generally, for a chosen target section $r:P\to H$ and current allocation
$a$, an available feedback-dependent lift
$h(p,a)=r(pa)r(a)^{-1}$ telescopes along histories without requiring $r$ to be
a homomorphism. Restricted actuators, feedback and task relevance determine
the cost. With only a final checkpoint the fixed-actuator cost is periodic,
not linear in $T$. A free destructive reset changes the operation set and may
erase task information. Conversely, a split extension can have positive
tracking cost under restricted actuators and chosen targets. None of these
costs is an intrinsic welfare interpretation of nonsplitting.

\section{Guaranteed success does not elicit private costs}
Consider $n$ agents with independent positive finite-support private costs
$C_i<B$. An active set $S$ succeeds exactly when $S\in\mathcal F$, where
$\mathcal F$ is a known nonempty upward-closed family and
$\varnothing\notin\mathcal F$. Success gives every agent benefit $B$;
agent $i$ pays $C_i$ iff active. There are no transfers, verification or
type-dependent action restrictions. A finite communication protocol may use
public and private randomization independent of costs. We consider an equilibrium
of the whole reporting, communication and action game, not obedience to a
mediator who already knows the costs.

\begin{proposition}[Guaranteed-success frontier]
Among such equilibria succeeding surely at every supported type profile,
maximum expected welfare is
\[
 nB-\min_{S\in\mathcal F}\sum_{i\in S}\mathbb E C_i.
\]
First-best welfare is attainable exactly when some fixed $S\in\mathcal F$
is realized-cost-minimizing almost surely.
\end{proposition}
\begin{proof}
Let $x_i(c)$ be agent $i$'s interim probability of activity. A type $c$ can
mimic another type $c'$'s entire randomized communication and action strategy.
Independence makes the opponents' type law the same, and guaranteed success
makes this deviation's success benefit $B$. Hence
$B-cx_i(c)\ge B-cx_i(c')$. Since $c>0$, interchanging types gives
$x_i(c)=x_i(c')=:x_i$. Thus expected cost is
$\sum_i\mathbb E[C_i]x_i$. The unconditional active-set law is supported on
$\mathcal F$, so this sum is at least the minimum displayed coalition cost.
Choose a coalition attaining that minimum. Positive means imply it is
inclusion-minimal: its members are pivotal, prefer $B-C_i>0$ to failure,
and outsiders prefer abstention. This designated-coalition equilibrium attains
the bound. Finally the minimum of expected coalition costs equals the expected
minimum realized cost iff a fixed minimizing coalition attains that realized
minimum almost surely, since their nonnegative difference must vanish.
\end{proof}

This is a finite-protocol specialization of established screening reasoning,
not a new general incentive principle. Battaglini--Palfrey's honest-and-obedient
constraints already include joint report/action deviations\citep{battaglini2024}.
The statement covers threshold and coverage tasks but requires independence,
guaranteed success and the stated common-benefit utility. It is weak
implementation of one equilibrium, not uniqueness or full implementation.

For a calibration, let $n=2$, $C_i\in\{l,h\}$ independently, with low
probability $p$, $0<l<h<B$, $a=1-p$ and $\bar c=pl+ah$.
The unconstrained team optimum over labelled local type policies and
type-independent public randomness is
\[
 W_{\rm PL}=\max\{2B-\bar c,\;2B(2p-p^2)-2pl\}.
\]
Indeed a multilinear objective attains its maximum at deterministic local
policies. Conditional on the other's policy, activity is worthwhile at cost
$c$ iff $c\le B\Pr(\text{other inactive})$; thus the high-only rule is
unnecessary, and checking the remaining vertices gives the formula.
An externally informed obedient mediator attains
$W_*=2B-l(2p-p^2)-ha^2$, with gap
\[
 W_*-W_{\rm PL}
 =\min\{pa(h-l),\;(2B-h)a^2+lp^2\}>0.
\]
Equality occurs only at $p\in\{0,1\}$ or $h=l$ in this closed parameter
domain. Neither this team bound nor the externally informed mediator is a
truthful cheap-talk implementation. With transfers conditional on verified
completion, second-price procurement admits an ex-post honest-and-obedient
equilibrium; arbitrary disobedient opponents prevent the stronger claim of
dominant strategies in the full voluntary-action game.

\section{Pricing trusted verification}
Now restrict the architecture: the coordinator receives no private reports or
other type-dependent signals, and can query an external truthful verifier for
one previously uninspected actual cost per unit. Histories, stopping and
selection depend only on those answers and independent randomness. After
inspection it designates one worker, who obeys since $B>C_i$.
For arbitrary $n\ge2$ iid binary costs as above, put
\[
 W_0=nB-\bar c,\quad v=pa(h-l),\quad
 R_*=\sum_{j=0}^{n-2}a^j=\frac{1-a^{n-1}}p.
\]
\begin{proposition}[Exact expected-query frontier]
For expected query budget $r\ge0$, maximum gross welfare is
$W(r)=W_0+v\min\{r,R_*\}$.
With a hard cap $m\le n-1$, it is $nB-l-(h-l)a^{m+1}$.
Zero-error first best requires worst-case $n-1$ queries and minimum expected
count $R_*$.
\end{proposition}
\begin{proof}
Before finding a low cost, with at least two uninspected agents, the best
fallback has expected cost $\bar c$. A query finds low with probability $p$
and reduces cost to $l$; otherwise an uninspected fallback remains. Each such
useful query has conditional gain $p(\bar c-l)=v$. After finding low, or with
only one uninspected agent, queries have zero value. If $U$ counts useful
queries, conditional expectation gives gain at most $v\mathbb E U$.
The $j$th useful query occurs with probability at most $a^{j-1}$, so
$\mathbb E U\le\min\{r,R_*\}$. Search until the first low or $n-1$ highs,
then select the low or uninspected fallback. This attains first best with
expected count $R_*$. Mixing this whole policy with no search attains every
intermediate expected-budget point. Capping search at $m$ gives expected
count $\sum_{j=0}^{m-1}a^j$ and the stated value. Finally, after fewer than
$n-1$ high observations, stopping can select a high while an uninspected low
exists with positive probability. Zero-error selection must continue on all
such histories, proving both lower bounds.
\end{proof}
A uniform random inspection order gives an equivariant law with unchanged
query count; full search selects uniformly among cheapest agents.
For per-query price $\kappa\ge0$, net optimum is
$W_0+\max\{0,v-\kappa\}R_*$. Query addresses, role labels and recommendation
bits are outside this resource count. This is a special nonobligatory-inspection
problem\citep{beyhaghi2022}, not a general minimum-communication theorem.

\section{Action-orbit algorithms and the input boundary}
Consider complete-information rational games with explicitly listed actions
and payoffs $u_i(s)=\sum_{F\in\mathcal F_i}f_{iF}(s_F)$, where $i\in F$,
$|F|\le d$ and each factor table is explicit. Let $L$ include all table,
scope, coefficient and supplied typed-group generator encodings. Exact
automorphism validity is a promise; its discovery is not free.
Let $q$ count group orbits of the disjoint player-action pairs $(i,a)$.
An invariant independent profile assigns probability $x_w$ to each action in
orbit $w$, with normalizations $\sum_w c_{iw}x_w=1$.

\begin{proposition}[Local-factor orbit reduction]
Conditional payoff polynomials $P_w(x)$ of degree at most $d-1$ can be built
in polynomial time in $L$. Invariant Nash equilibria are exactly solutions of
\[
 x_w\ge0,\quad \sum_w c_{iw}x_w=1,\quad
 P_w(x)\le v_{B(w)},\quad
 x_w(v_{B(w)}-P_w(x))=0,
\]
where $B(w)$ is the player orbit. There are at most $2q$ variables and
$O(q)$ constraints of degree at most $D=\max\{2,d\}$.
\end{proposition}
\begin{proof}
For a representative $(i,a)$, sum each factor table against the independent
probabilities of its other players, fixing $a$. This constructs $P_w$ with
at most $d-1$ probability factors per term. Exact typed automorphisms identify
conditional payoffs throughout the orbit. There are at most $q$ player orbits;
the displayed constraints say precisely that every action's payoff is at most
the player's value and every positive-probability action attains it. These
are the mixed Nash conditions. The classical invariant-equilibrium existence
argument applies, so the system is nonempty.
\end{proof}

Using standard real-algebraic sampling, this gives runtime
$(qD)^{O(q)}\operatorname{poly}(L)$ for an algebraically represented exact
solution. The coefficient-bit and output-encoding bound is an external
algorithmic ingredient, not implemented here; its precise original-source
formulation remains to be verified before publication. The primary algorithmic
baseline is Tewolde et al.'s action-orbit method\citep{tewolde2025}; symmetric
graphical-game algorithms are also prior art\citep{brandt2011}.

There is a useful input-size refinement. Remove singleton-action players and
substitute their actions. A residual factor on $t$ active players has at least
$2^t$ table entries, so $t\le\log_2 L$. Consequently the bound is
$(q\log L)^{O(q)}\operatorname{poly}(L)=f(q)\operatorname{poly}(L)$:
for $a=Cq$, $(\log L)^a\le(a/e)^aL$, up to fixed log-base factors.
The same observation applies to the explicit normal-form baseline; it does
not establish a new separation from that prior algorithm. Expanding the
orbit coordinates to individual actions costs at least the output length.
Rational approximation needs a normalization repair: approximate orbit
masses within each player's action-orbit blocks, renormalize their sum, and
divide by block multiplicities. Payoff Lipschitz bounds from the explicit
rational tables give an $\varepsilon$-Nash output at precision polynomial in
$L$ and $\log(1/\varepsilon)$.

One can also append a polynomial welfare threshold to this compact system
and bisect its feasible values, obtaining an $\eta$-optimal equilibrium
\emph{within the invariant Nash set}. Compactness gives attainment; rational
rounding requires an additional welfare-error allowance. This does not
optimize unrestricted Nash or guarantee good welfare.

An invariant correlated-equilibrium program instead uses probabilities on
\emph{profile} orbits and representative deviation-orbit inequalities.
Averaging preserves CE and an invariant linear welfare objective; an optimal
solution can use at most $r+1$ profile orbits when there are $r$ deviation rows.
However binary cyclic games can have at least $2^n/n$ profile orbits despite
only two action orbits. Neither enumeration nor decentralized correlated
sampling is free. Compact CE algorithms are established prior
work\citep{papadimitriou2008,jiang2011}.

\subsection{A one-orbit counting obstruction}
Given a Boolean circuit $F$ on $n\ge2$ bits, let $r$ rotate coordinates and
$c$ complement them. Give all players the same succinct payoff circuit
\[
 u_i(s)=\frac1{2n}\sum_{j=0}^{n-1}\bigl(F(r^js)+F(cr^js)\bigr).
\]
Rotation and simultaneous complementation are exact typed automorphisms,
transitive on player-action pairs: $q=1$. The unique invariant product profile
is fair independent bits. It is Nash because complementation equates each
player's two conditional payoffs. Yet
\[
 \mathbb E u_i=2^{-n}\#\{s:F(s)=1\}.
\]
Every transformation permutes the input cube, proving the equality and an
explicit counting-SAT reduction for exact value evaluation. This standard-style
reduction concerns succinct circuits, not explicit local tables. Finding the
uniform equilibrium is trivial here; invariant regret is zero. No hardness
of equilibrium output, regret evaluation or value-free near-optimal profile
selection follows.

\section{When symmetry averaging changes implementation}
Let a finite mechanism's feasible vector set be nonempty compact convex and
preserved by a finite relabelling group. Include linear joint report/action
incentives and a group-invariant linear resource budget. For an invariant
\emph{affine} objective $J$, averaging an optimizer over the group gives an
invariant optimizer at the same resource level. This is established reasoning,
explicit in collective-action mechanism design\citep{battaglini2024}.
Invariance alone is insufficient: on the two-coordinate probability simplex,
$J(z)=z_1^2+z_2^2$ has only noninvariant vertex maximizers under exchange.

Nor does averaging preserve product-law implementations. Independent actions
produce exactly one active agent with probability one only by designating a
surely active agent and making all others surely inactive. Their permutation
average is correlated, so public role selection changes the architecture.
At a fixed input, allowed outputs decompose into stabilizer orbits. An invariant
law is a mixture of uniform orbit laws, with minimum entropy $\log_2 m$ for
smallest orbit size $m$. A deterministic decoder of $b$ fair seed bits exists
exactly when $2^b$ is a nonnegative integer combination of allowed orbit sizes.
A single three-element orbit therefore has no fixed finite fair-bit sampler;
rejection using two bits has expected cost $8/3$ bits and unbounded worst case.
These are centralized sampling facts, not general distributed communication
or anonymous-election bounds.

\section{Incentives, learning and remaining questions}\label{sec:06-outlook-1}
Potential structure guarantees neither efficient selection nor suitable
information. The preserved feasibility example is fully symmetric and has
an exact potential whose global maximum is the worst-welfare pure equilibrium.
A separate three-player common-interest coordination game, paying
$\epsilon>0$ at all-zero, $1$ at all-one and zero otherwise, has precisely
two strict pure equilibria for $0<\epsilon<1$. Its pure-equilibrium price of
anarchy is $1/\epsilon$ and price of stability is $1$. This elementary control
shows no efficiency bound follows from full ordinary symmetry alone.
The claim does not extend unchanged to four players: balanced profiles add
weak zero-welfare equilibria.

Strict unilateral improvements terminate in finite potential games
\citep{monderer1996}; simultaneous best responses can cycle. Under uniform
asynchronous log-linear updates with weights $\lambda^{u_i}$, the finite
exact-potential model has stationary law proportional to $\lambda^\Phi$ by
detailed balance. Its high-$\lambda$ limit selects potential maximizers, not
necessarily welfare maximizers. Full assumptions, proof and an exact rational
chain check are kept in research; no general learning or LLM convergence
claim follows. The companion binary-potential candidate supplies another
possible source of potential existence, under its unchanged hypotheses.

These initial results support a coherent research programme, with no originality
clearance yet. The most promising next question adds decentralized private
communication, verification or incentive constraints to the exact Bayesian
architecture gap. A second lead optimizes genuine checkpoint tasks under
partial feedback and shared actuator costs. Information budgets must distinguish
scalar measurements, precision, samples, bits and computation. Neither a
standard averaging proof nor a small control example alone warrants a broad
new theory of agentic AI. Empirical language-model experiments remain separate
from the idealized agents analyzed here.

\section*{Acknowledgements and computational assistance}
OpenAI Codex assisted with source comparison, proof reconstruction, exact
computations and manuscript preparation. Internal AI reviews are not external
peer review. The author remains responsible for the mathematics and its
attribution. Source versions and computational checks are recorded in the
accompanying repository.
\begingroup\small
\interlinepenalty=10000
\bibliographystyle{plainnat}
\bibliography{references}
\endgroup
\end{document}
